





















































































































\begin{register}{H}{INTERRUPT\_CORE0\_\regindex{GPIO\_INTERRUPT\_PRO\_MAP}\_REG}{0x0034}\label{regdesc:INTERRUPTCORE0GPIOINTERRUPTPROMAPREG}






\end{register}\vspace{-3em}


















\begin{register}{H}{INTERRUPT\_CORE0\_\regindex{SPI\_INTR\_2\_MAP}\_REG}{0x0040}\label{regdesc:INTERRUPTCORE0SPIINTR2MAPREG}






\end{register}\vspace{-3em}










\begin{register}{H}{INTERRUPT\_CORE0\_\regindex{UART\_INTR\_MAP}\_REG}{0x0044}\label{regdesc:INTERRUPTCORE0UARTINTRMAPREG}






\end{register}\vspace{-3em}


\begin{register}{H}{INTERRUPT\_CORE0\_\regindex{UART1\_INTR\_MAP}\_REG}{0x0048}\label{regdesc:INTERRUPTCORE0UART1INTRMAPREG}






\end{register}\vspace{-3em}


\begin{register}{H}{INTERRUPT\_CORE0\_\regindex{LEDC\_INT\_MAP}\_REG}{0x004C}\label{regdesc:INTERRUPTCORE0LEDCINTMAPREG}






\end{register}\vspace{-3em}


\begin{register}{H}{INTERRUPT\_CORE0\_\regindex{EFUSE\_INT\_MAP}\_REG}{0x0050}\label{regdesc:INTERRUPTCORE0EFUSEINTMAPREG}






\end{register}\vspace{-3em}


















\begin{register}{H}{INTERRUPT\_CORE0\_\regindex{RTC\_CORE\_INTR\_MAP}\_REG}{0x0054}\label{regdesc:INTERRUPTCORE0RTCCOREINTRMAPREG}






\end{register}\vspace{-3em}










\begin{register}{H}{INTERRUPT\_CORE0\_\regindex{I2C\_EXT0\_INTR\_MAP}\_REG}{0x058}\label{regdesc:INTERRUPTCORE0I2CEXT0INTRMAPREG}






\end{register}\vspace{-3em}


















\begin{register}{H}{INTERRUPT\_CORE0\_\regindex{TG\_T0\_INT\_MAP}\_REG}{0x005C}\label{regdesc:INTERRUPTCORE0TGT0INTMAPREG}






\end{register}\vspace{-3em}


\begin{register}{H}{INTERRUPT\_CORE0\_\regindex{TG\_WDT\_INT\_MAP}\_REG}{0x0060}\label{regdesc:INTERRUPTCORE0TGWDTINTMAPREG}






\end{register}\vspace{-3em}


























\begin{register}{H}{INTERRUPT\_CORE0\_\regindex{SYSTIMER\_TARGET0\_INT\_MAP}\_REG}{0x0068}\label{regdesc:INTERRUPTCORE0SYSTIMERTARGET0INTMAPREG}






\end{register}\vspace{-3em}


\begin{register}{H}{INTERRUPT\_CORE0\_\regindex{SYSTIMER\_TARGET1\_INT\_MAP}\_REG}{0x006C}\label{regdesc:INTERRUPTCORE0SYSTIMERTARGET1INTMAPREG}






\end{register}\vspace{-3em}


\begin{register}{H}{INTERRUPT\_CORE0\_\regindex{SYSTIMER\_TARGET2\_INT\_MAP}\_REG}{0x0070}\label{regdesc:INTERRUPTCORE0SYSTIMERTARGET2INTMAPREG}






\end{register}\vspace{-3em}


























\begin{register}{H}{INTERRUPT\_CORE0\_\regindex{APB\_ADC\_INT\_MAP}\_REG}{0x0080}\label{regdesc:INTERRUPTCORE0APBADCINTMAPREG}






\end{register}\vspace{-3em}


\begin{register}{H}{INTERRUPT\_CORE0\_\regindex{DMA\_CH0\_INT\_MAP}\_REG}{0x0084}\label{regdesc:INTERRUPTCORE0DMACH0INTMAPREG}






\end{register}\vspace{-3em}


































\begin{register}{H}{INTERRUPT\_CORE0\_\regindex{SHA\_INT\_MAP}\_REG}{0x0088}\label{regdesc:INTERRUPTCORE0SHAINTMAPREG}






\end{register}\vspace{-3em}


\begin{register}{H}{INTERRUPT\_CORE0\_\regindex{ECC\_INT\_MAP}\_REG}{0x008C}\label{regdesc:INTERRUPTCORE0ECCINTMAPREG}






\end{register}\vspace{-3em}


\begin{register}{H}{INTERRUPT\_CORE0\_\regindex{CPU\_INTR\_FROM\_CPU\_0\_MAP}\_REG}{0x0090}\label{regdesc:INTERRUPTCORE0CPUINTRFROMCPU0MAPREG}






\end{register}\vspace{-3em}


\begin{register}{H}{INTERRUPT\_CORE0\_\regindex{CPU\_INTR\_FROM\_CPU\_1\_MAP}\_REG}{0x0094}\label{regdesc:INTERRUPTCORE0CPUINTRFROMCPU1MAPREG}






\end{register}\vspace{-3em}


\begin{register}{H}{INTERRUPT\_CORE0\_\regindex{CPU\_INTR\_FROM\_CPU\_2\_MAP}\_REG}{0x0098}\label{regdesc:INTERRUPTCORE0CPUINTRFROMCPU2MAPREG}






\end{register}\vspace{-3em}


\begin{register}{H}{INTERRUPT\_CORE0\_\regindex{CPU\_INTR\_FROM\_CPU\_3\_MAP}\_REG}{0x009C}\label{regdesc:INTERRUPTCORE0CPUINTRFROMCPU3MAPREG}






\end{register}\vspace{-3em}


\begin{register}{H}{INTERRUPT\_CORE0\_\regindex{ASSIST\_DEBUG\_INTR\_MAP}\_REG}{0x00A0}\label{regdesc:INTERRUPTCORE0ASSISTDEBUGINTRMAPREG}






\end{register}\vspace{-3em}

















































\begin{register}{H}{INTERRUPT\_CORE0\_\regindex{PIF\_PMS\_MONITOR\_VIOLATE\_SIZE\_INTR\_MAP}\_REG}{0x00A4}\label{regdesc:INTERRUPTCORE0PIFPMSMONITORVIOLATESIZEINTRMAPREG}
\regfield{(reserved)}{27}{5}{000000000000000000000000000}%
\regfield{INTERRUPT\_CORE0\_SOURCE\_\regindex{X}\_MAP}{5}{0}{{0}}%
\reglabel{Reset}\regnewline%
\begin{regdesc}\begin{reglist}
\label{fielddesc:INTERRUPTCORE0CACHECORE0ACSINTMAP}\item [INTERRUPT\_CORE0\_SOURCE\_\regindex{X}\_MAP] 将中断源 SOURCE\_\regindex{X} 映射至 CPU 外部中断。中断源 SOURCE\_\regindex{X} 见表 \ref{tab:peripheral-interrupt-config-status-source}。(R/W)
\end{reglist}\end{regdesc}
\end{register}


\begin{register}{H}{INTERRUPT\_CORE0\_INTR\_STATUS\_0\_REG}{0x00AC}\label{regdesc:INTERRUPTCORE0INTRSTATUS0REG}
\regfield{INTERRUPT\_CORE0\_INTR\_STATUS\_0}{32}{0}{{0x000000}}%
\reglabel{Reset}\regnewline%
\begin{regdesc}\begin{reglist}
\label{fielddesc:INTERRUPTCORE0INTRSTATUS0}\item [INTERRUPT\_CORE0\_INTR\_STATUS\_0] 用于存储外部中断源的状态，每一位均代表一个外部中断源的状态，对应中断编号源：0 \verb+~+ 31。如果对应的位为 1，则表示该中断源触发了中断。(RO)
\end{reglist}\end{regdesc}
\end{register}


\begin{register}{H}{INTERRUPT\_CORE0\_INTR\_STATUS\_1\_REG}{0x00B0}\label{regdesc:INTERRUPTCORE0INTRSTATUS1REG}
\regfield{INTERRUPT\_CORE0\_INTR\_STATUS\_1}{32}{0}{{0x000000}}%
\reglabel{Reset}\regnewline%
\begin{regdesc}\begin{reglist}
\label{fielddesc:INTERRUPTCORE0INTRSTATUS1}\item [INTERRUPT\_CORE0\_INTR\_STATUS\_1] 用于存储外部中断源的状态，每一位均代表一个外部中断源的状态，对应中断编号源：32 \verb+~+ 42。如果对应的位为 1，则表示该中断源触发了中断。(RO)
\end{reglist}\end{regdesc}
\end{register}


\begin{register}{H}{INTERRUPT\_CORE0\_CLOCK\_GATE\_REG}{0x00B4}\label{regdesc:INTERRUPTCORE0CLOCKGATEREG}
\regfield{(reserved)}{31}{1}{0000000000000000000000000000000}%
\regfield{INTERRUPT\_CORE0\_CLK\_EN}{1}{0}{{1}}%
\reglabel{Reset}\regnewline%
\begin{regdesc}\begin{reglist}
\label{fielddesc:INTERRUPTCORE0CLKEN}\item [INTERRUPT\_CORE0\_CLK\_EN] 置 1 强制使能中断寄存器的时钟门控。(R/W)

\end{reglist}\end{regdesc}
\end{register}


\begin{register}{H}{INTERRUPT\_CORE0\_CPU\_INT\_ENABLE\_REG}{0x00B8}\label{regdesc:INTERRUPTCORE0CPUINTENABLEREG}
\regfield{INTERRUPT\_CORE0\_CPU\_INT\_ENABLE}{32}{0}{{0}}%
\reglabel{Reset}\regnewline%
\begin{regdesc}\begin{reglist}
\label{fielddesc:INTERRUPTCORE0CPUINTENABLE}\item [INTERRUPT\_CORE0\_CPU\_INT\_ENABLE] 在相应位写 1，即可使能对应 CPU 中断。更多配置信息，见章节 \ref{mod:riscvcpu} \textit{\nameref{mod:riscvcpu}}。(R/W)
\end{reglist}\end{regdesc}
\end{register}


\begin{register}{H}{INTERRUPT\_CORE0\_CPU\_INT\_TYPE\_REG}{0x00BC}\label{regdesc:INTERRUPTCORE0CPUINTTYPEREG}
\regfield{INTERRUPT\_CORE0\_CPU\_INT\_TYPE}{32}{0}{{0}}%
\reglabel{Reset}\regnewline%
\begin{regdesc}\begin{reglist}
\label{fielddesc:INTERRUPTCORE0CPUINTTYPE}\item [INTERRUPT\_CORE0\_CPU\_INT\_TYPE] 配置 CPU 中断类型。0：电平触发；1：边沿触发。更多配置信息，见章节 \ref{mod:riscvcpu} \textit{\nameref{mod:riscvcpu}}。(R/W)
\end{reglist}\end{regdesc}
\end{register}


\begin{register}{H}{INTERRUPT\_CORE0\_CPU\_INT\_CLEAR\_REG}{0x00C0}\label{regdesc:INTERRUPTCORE0CPUINTCLEARREG}
\regfield{INTERRUPT\_CORE0\_CPU\_INT\_CLEAR}{32}{0}{{0}}%
\reglabel{Reset}\regnewline%
\begin{regdesc}\begin{reglist}
\label{fielddesc:INTERRUPTCORE0CPUINTCLEAR}\item [INTERRUPT\_CORE0\_CPU\_INT\_CLEAR] 在相应位写 1，即可清除对应的 CPU 中断。更多配置信息，见章节 \ref{mod:riscvcpu} \textit{\nameref{mod:riscvcpu}}。(R/W)
\end{reglist}\end{regdesc}
\end{register}


\begin{register}{H}{INTERRUPT\_CORE0\_CPU\_INT\_EIP\_STATUS\_REG}{0x00C4}\label{regdesc:INTERRUPTCORE0CPUINTEIPSTATUSREG}
\regfield{INTERRUPT\_CORE0\_CPU\_INT\_EIP\_STATUS}{32}{0}{{0}}%
\reglabel{Reset}\regnewline%
\begin{regdesc}\begin{reglist}
\label{fielddesc:INTERRUPTCORE0CPUINTEIPSTATUS}\item [INTERRUPT\_CORE0\_CPU\_INT\_EIP\_STATUS] 用于存储 CPU 中断的阻塞状态。更多信息，请参考章节 \ref{mod:riscvcpu} \textit{\nameref{mod:riscvcpu}}。(RO)
\end{reglist}\end{regdesc}
\end{register}


\begin{register}{H}{INTERRUPT\_CORE0\_CPU\_INT\_PRI\_\regindex{n}\_REG}{\regindex{n}: 1 - 31)(0x00C8 + 0x4*\regindex{n}}\label{regdesc:INTERRUPTCORE0CPUINTPRI0REG}
\regfield{(reserved)}{28}{4}{0000000000000000000000000000}%
\regfield{INTERRUPT\_CORE0\_CPU\_PRI\_\regindex{n}\_MAP}{4}{0}{{0}}%
\reglabel{Reset}\regnewline%
\begin{regdesc}\begin{reglist}
\label{fielddesc:INTERRUPTCORE0CPUPRI0MAP}\item [INTERRUPT\_CORE0\_CPU\_PRI\_\regindex{n}\_MAP] 用于设置 CPU 中断 \regindex{n} 的优先级，可配置为 1 \verb+~+ 15。数字越大，优先级越高。更多配置信息，见章节 \ref{mod:riscvcpu} \textit{\nameref{mod:riscvcpu}}。(R/W)
\end{reglist}\end{regdesc}
\end{register}

\begin{register}{H}{INTERRUPT\_CORE0\_CPU\_INT\_THRESH\_REG}{0x0148}\label{regdesc:INTERRUPTCORE0CPUINTTHRESHREG}
\regfield{(reserved)}{28}{4}{0000000000000000000000000000}%
\regfield{INTERRUPT\_CORE0\_CPU\_INT\_THRESH}{4}{0}{{0}}%
\reglabel{Reset}\regnewline%
\begin{regdesc}\begin{reglist}
\label{fielddesc:INTERRUPTCORE0CPUINTTHRESH}\item [INTERRUPT\_CORE0\_CPU\_INT\_THRESH] 用于设置 CPU 中断阈值。仅当中断的优先级等于或高于该阈值，CPU 才会响应该中断。更多配置信息，见章节 \ref{mod:riscvcpu} \textit{\nameref{mod:riscvcpu}}。(R/W)

\end{reglist}\end{regdesc}
\end{register}


\begin{register}{H}{INTERRUPT\_CORE0\_INTERRUPT\_DATE\_REG}{0x07FC}\label{regdesc:INTERRUPTCORE0INTERRUPTDATEREG}
\regfield{(reserved)}{4}{28}{0000}%
\regfield{INTERRUPT\_CORE0\_INTERRUPT\_DATE}{28}{0}{{0x2108190}}%
\reglabel{Reset}\regnewline%
\begin{regdesc}\begin{reglist}
\label{fielddesc:INTERRUPTCORE0INTERRUPTDATE}\item [INTERRUPT\_CORE0\_INTERRUPT\_DATE] 版本寄存器。(R/W)

\end{reglist}\end{regdesc}
\end{register}
